\documentclass[../../main.tex]{subfiles}% 注意这里的文件路径不能用 ./main.tex ,否则用latexmk编译子文件会报错
\graphicspath{{\subfix{./image/}}{\subfix{../../image/}}} % 指定图片引用路径,后续可以直接使用图片文件名
% 如果图片存放在上级目录,则使用 ../../image/ 作为备用路径

% 例如:
% \begin{figure}[H]
% \centering
% \includegraphics[width=0.3\textwidth]{图.png}
% \caption{}
% \label{figure:图}
% \end{figure}
% 注意:上述\label{}一定要放在\caption{}之后,否则引用图片序号会只会显示??.

\begin{document}

\section{基本解和 Green 函数}

\subsection{基本解}

\begin{definition}
对 $x\in\mathbb{R}^n$，$x\neq 0$，称函数
\begin{align}
\Gamma(x) = \begin{cases}
-\frac{1}{2\pi}\ln|x|, & n=2, \\
\frac{1}{n(n-2)\alpha(n)|x|^{n-2}}, & n\geqslant 3
\end{cases} \label{eq:fundamental}
\end{align}
为\textbf{Laplace 方程的基本解}。
\end{definition}
\begin{remark}
Laplace方程的基本解具有很强的物理意义：当 $n=3$ 时，基本解 $\Gamma(x)$ 就是由放置在原点的单位正点电荷引起的在全空间 $\mathbb{R}^3$ 上的静电位势分布。
\end{remark}

\begin{proposition}
设$\Gamma (x)$是Laplace方程的基本解.对于 $x\neq 0$，有
\begin{align*}
|D\Gamma(x)| \leqslant \frac{C}{|x|^{n-1}},\qquad |D^2\Gamma(x)| \leqslant \frac{C}{|x|^n},
\end{align*}
其中 $C$ 是一个仅依赖于空间维数 $n$ 的正常数。
\end{proposition}
\begin{remark}
我们通常用 $L^1(\Omega)$ 表示 $\Omega$ 上的可积函数空间，用 $L^1_{\mathrm{loc}}(\Omega)$ 表示 $\Omega$ 上的局部可积函数空间。实际上，基本解 $\Gamma(x)\in L^1_{\mathrm{loc}}(\mathbb{R}^n)$ 且在广义函数意义下满足方程
\begin{align*}
-\Delta \Gamma(x) = \delta(x),
\end{align*}
其中 $\delta(x)$ 是 Dirac 测度。也就是说，对于任意 $f\in C_0^\infty(\mathbb{R}^n)$，成立
\begin{align*}
\int_{\mathbb{R}^n} \Gamma(x)[-\Delta f(x)]\,\mathrm{d}x = f(0).
\end{align*}
\end{remark}

注意到当 $x\neq 0$ 时，$\Gamma(x)$ 满足 Laplace 方程 。作平移变换，我们知道，当 $x\neq y$ 时，$\Gamma(x-y)$ 满足 Laplace 方程。由于方程   是线性的，我们尝试求下列形式的解：
\begin{align*}
u(x) = \int_{\mathbb{R}^n} \Gamma(x-y)f(y)\,\mathrm{d}y.
\end{align*}
然而，由于 $D^2\Gamma(x-y)$ 在 $x=y$ 附近不可积，故 $u(x)$ 并不满足 Laplace 方程  。

\begin{lemma}
如果 $u,v\in C^2(\Omega)\cap C^1(\overline{\Omega})$，则
\begin{align}
\int_\Omega (u\Delta v - v\Delta u)\,\mathrm{d}x = \int_{\partial\Omega} \left( u\frac{\partial v}{\partial n} - v\frac{\partial u}{\partial n} \right) \mathrm{d}S(x), \label{eq:green-second}
\end{align}
其中 $n$ 是 $\partial\Omega$ 的单位外法向量。
\end{lemma}

\begin{theorem}\label{theorem:Laplace方程基本解的性质}
假设 $f(x)\in C_0^2(\mathbb{R}^n)$，定义
\begin{align*}
u(x) = \int_{\mathbb{R}^n} \Gamma(x-y)f(y)\,\mathrm{d}y,
\end{align*}
其中$\Gamma(x)$是Laplace方程的基本解.则 $u\in C^2(\mathbb{R}^n)$ 且满足 $-\Delta u=f$。
\end{theorem}
\begin{remark}
此定理说明我们可以利用基本解来构造位势方程(Poisson 方程)在全空间上的解。
\end{remark}
\begin{proof}
固定 $x\in\mathbb{R}^n$。

(1) 我们注意到
\begin{align*}
u(x) = \int_{\mathbb{R}^n} \Gamma(y)f(x-y)\,\mathrm{d}y,
\end{align*}
从而有
\begin{align*}
\frac{u(x+he_i)-u(x)}{h} = \int_{\mathbb{R}^n} \Gamma(y)\left[ \frac{f(x+he_i-y)-f(x-y)}{h} \right] \mathrm{d}y,
\end{align*}
其中 $e_i$ 表示 $\mathbb{R}^n$ 上的第 $i$ 个方向向量。由于 $f$ 具有紧支集，则
\begin{align*}
\lim_{h\to 0} \frac{f(x+he_i-y)-f(x-y)}{h} = \frac{\partial}{\partial x_i} f(x-y),
\end{align*}
且上述极限关于 $y$ 一致收敛。于是
\begin{align*}
\frac{\partial u}{\partial x_i} = \int_{\mathbb{R}^n} \Gamma(y) \frac{\partial}{\partial x_i} f(x-y)\,\mathrm{d}y \quad (i=1,2,\cdots,n).
\end{align*}
同理我们得到
\begin{align*}
\frac{\partial^2 u}{\partial x_i\partial x_j} = \int_{\mathbb{R}^n} \Gamma(y) \frac{\partial^2}{\partial x_i\partial x_j} f(x-y)\,\mathrm{d}y \quad (i,j=1,2,\cdots,n).
\end{align*}
由于上式右端项对变量 $x$ 是连续的，因而 $u\in C^2(\mathbb{R}^n)$。实际上，由于 $f$ 具有紧支集，我们有：存在 $M_0>0$，使得对所有 $x\in\mathbb{R}^n$，有
\begin{align*}
|f(x)| + |Df(x)| + |D^2f(x)| \leqslant M_0.
\end{align*}
由此容易证明 $u, Du, D^2u$ 是有界的，即存在 $M_1>0$ 使得对所有 $x\in\mathbb{R}^n$，有
\begin{align*}
|u(x)| + |Du(x)| + |D^2u(x)| \leqslant M_1.
\end{align*}

(2) 固定 $\varepsilon>0$，则
\begin{align*}
\Delta u &= \int_{\mathbb{R}^n} \Gamma(y)\Delta_x f(x-y)\,\mathrm{d}y = \int_{B(0,\varepsilon)} \Gamma(y)\Delta_x f(x-y)\,\mathrm{d}y + \int_{\mathbb{R}^n\setminus B(0,\varepsilon)} \Gamma(y)\Delta_x f(x-y)\,\mathrm{d}y,
\end{align*}
这里 $\Delta_x$ 是关于变量 $x$ 的 Laplace 算子 $\Delta$。记
\begin{align*}
I_\varepsilon = \int_{B(0,\varepsilon)} \Gamma(y)\Delta_x f(x-y)\,\mathrm{d}y, \quad J_\varepsilon = \int_{\mathbb{R}^n\setminus B(0,\varepsilon)} \Gamma(y)\Delta_x f(x-y)\,\mathrm{d}y,
\end{align*}
于是
\begin{align*}
\Delta u = I_\varepsilon + J_\varepsilon.
\end{align*}

我们首先估计 $I_\varepsilon$。利用函数 $f$ 的有界性我们容易证明
\begin{align}
|I_\varepsilon| \leqslant M_0 \int_{B(0,\varepsilon)} |\Gamma(y)|\,\mathrm{d}y
\leqslant \begin{cases}
C\varepsilon^2|\ln\varepsilon|, & n=2, \\
C\varepsilon^2, & n\geqslant 3,
\end{cases} \label{eq:Ieps}
\end{align}
其中 $C$ 是不依赖于 $\varepsilon$ 的常数。

由于函数 $f(y)\in C_0^2(\mathbb{R}^n)$ 具有紧支集，因而函数 $f(x-y)$ 具有紧支集。故存在一个正数 $R_x>0$，使得 $f(x-y)$ 的支集 $\operatorname{spt} f(x-y)\subset B(0,R_x)$。由引理 \eqref{eq:green-second} 我们可以将 $J_\varepsilon$ 表示成
\begin{align*}
J_\varepsilon &= \int_{\mathbb{R}^n\setminus B(0,\varepsilon)} \Gamma(y)\Delta_y f(x-y)\,\mathrm{d}y 
= \int_{B(0,R_x)\setminus B(0,\varepsilon)} \Gamma(y)\Delta_y f(x-y)\,\mathrm{d}y \\
&= \int_{B(0,R_x)\setminus B(0,\varepsilon)} \Delta\Gamma(y) f(x-y)\,\mathrm{d}y + \int_{\partial B(0,R_x)\cup\partial B(0,\varepsilon)} \left[ \Gamma(y)\frac{\partial}{\partial n} f(x-y) - \frac{\partial\Gamma(y)}{\partial n} f(x-y) \right] \mathrm{d}S(y) \\
&= \int_{\partial B(0,\varepsilon)} \left[ \Gamma(y)\frac{\partial}{\partial n} f(x-y) - \frac{\partial\Gamma(y)}{\partial n} f(x-y) \right] \mathrm{d}S(y),
\end{align*}
其中在球面 $\partial B(0,\varepsilon)$ 上 $n$ 表示单位内法向量，在球面 $\partial B(0,R_x)$ 上 $n$ 表示单位外法向量。记
\begin{align*}
K_\varepsilon = \int_{\partial B(0,\varepsilon)} \Gamma(y)\frac{\partial}{\partial n} f(x-y)\,\mathrm{d}S(y), \quad 
L_\varepsilon = -\int_{\partial B(0,\varepsilon)} \frac{\partial\Gamma(y)}{\partial n} f(x-y)\,\mathrm{d}S(y),
\end{align*}
于是
\begin{align*}
J_\varepsilon = K_\varepsilon + L_\varepsilon.
\end{align*}

我们首先估计 $K_\varepsilon$，然后再计算 $L_\varepsilon$。利用函数 $f$ 的梯度 $Df$ 的有界性我们容易证明
\begin{align}
|K_\varepsilon| \leqslant M_0 \int_{\partial B(0,\varepsilon)} |\Gamma(y)|\,\mathrm{d}S(y)
\leqslant \begin{cases}
C\varepsilon|\ln\varepsilon|, & n=2, \\
C\varepsilon, & n\geqslant 3,
\end{cases} \label{eq:Keps}
\end{align}
其中 $C$ 是不依赖于 $\varepsilon$ 的常数。

注意到在 $\partial B(0,\varepsilon)$ 上，有
\begin{align*}
D\Gamma(y) = \frac{-y}{n\alpha(n)|y|^n}, \qquad n = -\frac{y}{|y|} = -\frac{y}{\varepsilon},
\end{align*}
于是
\begin{align*}
\frac{\partial\Gamma(y)}{\partial n} = D\Gamma(y)\cdot n = \frac{1}{n\alpha(n)\varepsilon^{n-1}}.
\end{align*}
由中值公式我们得到
\begin{align}
L_\varepsilon = -\frac{1}{n\alpha(n)\varepsilon^{n-1}} \int_{\partial B(0,\varepsilon)} f(x-y)\,\mathrm{d}S(y)
= -\frac{f(x-y_\varepsilon)}{n\alpha(n)\varepsilon^{n-1}} \int_{\partial B(0,\varepsilon)} \mathrm{d}S(y) = -f(x-y_\varepsilon), \label{eq:Leps}
\end{align}
其中 $y_\varepsilon\in\partial B(0,\varepsilon)$。由函数 $f$ 的连续性，当 $\varepsilon\to 0$ 时，$L_\varepsilon\to -f(x)$。

综上我们知道，对任意 $\varepsilon>0$，有
\begin{align*}
\Delta u = I_\varepsilon + K_\varepsilon + L_\varepsilon.
\end{align*}
利用估计 \eqref{eq:Ieps}, \eqref{eq:Keps} 和 \eqref{eq:Leps}，令 $\varepsilon\to 0$，则有
\begin{align*}
-\Delta u(x) = f(x).
\end{align*}
至此我们完成定理的证明。
\end{proof}

\begin{theorem}
假设 $f(x)\in C_0^2(\mathbb{R}^n)$，则
\begin{align*}
u(x) = \int_{\mathbb{R}^n} \Gamma(x-y)f(y)\,\mathrm{d}y + C
\end{align*}
是位势方程在全空间 $\mathbb{R}^n$ 上所有的有界解。
\end{theorem}
\begin{proof}
利用\hyperref[theorem:Laplace方程基本解的性质]{Laplace方程基本解的性质}和 \hyperref[theorem:PDE-Liouville定理]{Liouville定理}得证。
\end{proof}

\begin{definition}
齐次边值问题
\begin{align}
\begin{cases}
-\Delta u = \lambda u, & x\in\Omega, \\
\left.\frac{\partial u}{\partial n}\right|_{\partial\Omega} = 0
\end{cases} \label{eq:eigenvalue}
\end{align}
称为\textbf{特征值问题}，使此问题有非零解的 $\lambda\in\mathbb{R}$ 称为此问题的\textbf{特征值}，相应的非零解称为对应于这个特征值的\textbf{特征函数}，记为 $u_\lambda$。
\end{definition}
\begin{remark}
在特征值问题 \eqref{eq:eigenvalue} 的方程两端同乘 $u$ 并在 $\Omega$ 上积分，我们不难证明特征值 $\lambda\geqslant 0$。
\end{remark}

\begin{theorem}
第二边值问题
\begin{align}
\begin{cases}
-\Delta u = \lambda u + f, & x\in\Omega, \\
\left.\frac{\partial u}{\partial n}\right|_{\partial\Omega} = g
\end{cases} \label{eq:neumann-gen}
\end{align}
有 $C^1(\overline{\Omega})\cap C^2(\Omega)$ 上的解的必要条件是对于相应的任意特征函数 $u_\lambda(x)$，成立
\begin{align}
\int_\Omega f(x)u_\lambda(x)\,\mathrm{d}x + \int_{\partial\Omega} g(x)u_\lambda(x)\,\mathrm{d}S(x) = 0. \label{eq:neumann-cond}
\end{align}
\end{theorem}

\begin{proof}
在引理 \eqref{eq:green-second} 公式中，取 $u$ 为第二边值问题 \eqref{eq:neumann-gen} 的解，$v=u_\lambda(x)$，立即得到要证的结论。
\end{proof}

\begin{corollary}
Poisson 方程的 Neumann 边值问题
\begin{align}
\begin{cases}
-\Delta u = f, & x\in\Omega, \\
\left.\frac{\partial u}{\partial n}\right|_{\partial\Omega} = g
\end{cases} \label{eq:neumann}
\end{align}
有 $C^1(\overline{\Omega})\cap C^2(\Omega)$ 上的解的必要条件是
\begin{align}
\int_\Omega f(x)\,\mathrm{d}x + \int_{\partial\Omega} g(x)\,\mathrm{d}S(x) = 0. \label{eq:neumann-cond0}
\end{align}
\end{corollary}
\begin{proof}
实际上，对于 Neumann 边值问题 \eqref{eq:neumann} 的方程两端在 $\Omega$ 上积分，利用 Gauss-Green 公式也可以得到上式。
\end{proof}


\subsection{Green 函数}

在这一小节中我们旨在得到 Dirichlet 问题  的求解公式。为此我们利用基本解来构造 Green 函数，并以此来获得解的表达式。

假设 $\Omega\subset\mathbb{R}^n$ 是有界开集且 $\partial\Omega$ 光滑，且 $u\in C^2(\Omega)\cap C^1(\overline{\Omega})$ 是 Dirichlet 问题  的解。固定 $x\in\Omega$。取充分小 $\varepsilon>0$，使得 $B(x,\varepsilon)\subset\Omega$。在区域 $\Omega\setminus B(x,\varepsilon)$ 上对 $u(y)$ 和基本解 $\Gamma(y-x)$ 应用引理 \eqref{eq:green-second} 的公式得
\begin{align*}
\int_{\Omega\setminus B(x,\varepsilon)} \left[ u(y)\Delta\Gamma(y-x) - \Gamma(y-x)\Delta u(y) \right] \mathrm{d}y
&= \int_{\partial\Omega} \left[ u(y)\frac{\partial}{\partial n}\Gamma(y-x) - \Gamma(y-x)\frac{\partial}{\partial n}u(y) \right] \mathrm{d}S(y) \\
&\quad + \int_{\partial B(x,\varepsilon)} \left[ u(y)\frac{\partial}{\partial n}\Gamma(y-x) - \Gamma(y-x)\frac{\partial}{\partial n}u(y) \right] \mathrm{d}S(y),
\end{align*}
其中在球面 $\partial B(x,\varepsilon)$ 上 $n$ 表示单位内法向量，在边界 $\partial\Omega$ 上 $n$ 表示单位外法向量。注意到，当 $x\neq y$ 时，$\Delta\Gamma(x-y)=0$，且当 $\varepsilon\to 0$ 时，
\begin{align*}
\int_{\partial B(x,\varepsilon)} u(y)\frac{\partial}{\partial n}\Gamma(y-x)\,\mathrm{d}S(y) \longrightarrow u(x),
\end{align*}
和
\begin{align*}
\int_{\partial B(x,\varepsilon)} \Gamma(y-x)\frac{\partial}{\partial n}u(y)\,\mathrm{d}S(y) \longrightarrow 0.
\end{align*}
从而得到
\begin{align}
u(x) = \int_{\partial\Omega} \left[ \Gamma(y-x)\frac{\partial}{\partial n}u(y) - u(y)\frac{\partial}{\partial n}\Gamma(y-x) \right] \mathrm{d}S(y) - \int_\Omega \Gamma(y-x)\Delta u(y)\,\mathrm{d}y. \label{eq:repr-integral}
\end{align}
但是，对于 Dirichlet 问题来说，$\frac{\partial}{\partial n}u(y)$ 仍然未知。我们将通过引进一个调和函数 $\phi^x$ 来消掉这一项。

假定对给定 $x\in\Omega$，函数 $\phi^x$ 满足方程
\begin{align}
\begin{cases}
-\Delta\phi^x(y) = 0, & y\in\Omega, \\
\phi^x|_{\partial\Omega} = \Gamma(y-x).
\end{cases} \label{eq:phi}
\end{align}
在 $\Omega$ 上对 $u(y)$ 和解 $\phi^x(y)$ 应用引理 \eqref{eq:green-second} 的公式得
\begin{align}
0 = \int_{\partial\Omega} \left[ u(y)\frac{\partial\phi^x}{\partial n} - \Gamma(y-x)\frac{\partial u}{\partial n} \right] \mathrm{d}S(y) + \int_\Omega \phi^x(y)\Delta u(y)\,\mathrm{d}y. \label{eq:phi-integral}
\end{align}
现在将 \eqref{eq:repr-integral} 和 \eqref{eq:phi-integral} 相加，于是我们得到
\begin{align*}
u(x) = -\int_{\partial\Omega} u(y)\frac{\partial}{\partial n}\left[ \Gamma(y-x) - \phi^x(y) \right] \mathrm{d}S(y) - \int_\Omega \left[ \Gamma(y-x) - \phi^x(y) \right] \Delta u(y)\,\mathrm{d}y.
\end{align*}

\begin{definition}
对于任意 $x,y\in\Omega$，$x\neq y$，函数
\begin{align*}
G(x,y) = \Gamma(y-x) - \phi^x(y)
\end{align*}
称为 $\Omega$ 上的 \textbf{Green 函数}。
\end{definition}

\begin{remark}
显然，当 $x\in\Omega$，$y\in\partial\Omega$ 时，$G(x,y)=0$。实际上，Green 函数就是基本解减去一个以基本解为边值的调和函数。由于函数 $\Gamma(y-x)$ 在 $\partial\Omega$ 附近是一个 $C^\infty$ 的光滑函数，当 $\Omega$ 满足一定的条件时，$\phi^x(y)$ 的存在性可以比较容易地得到（参阅文献[4]中第二章第八节）。
\end{remark}

据定义，上述 $u(x)$ 可表示为
\begin{align*}
u(x) = -\int_{\partial\Omega} u(y)\frac{\partial}{\partial n}G(x,y)\,\mathrm{d}S(y) - \int_\Omega G(x,y)\Delta u(y)\,\mathrm{d}y,
\end{align*}
于是我们得到下面的结论。

\begin{theorem}[解的表达式]
假设 $\Omega$ 是 $\mathbb{R}^n$ 上的一个有界区域，$u(x)\in C^2(\Omega)\cap C^1(\overline{\Omega})$ 是 Dirichlet 问题  的解，则
\begin{align}
u(x) = -\int_{\partial\Omega} g(y)\frac{\partial}{\partial n}G(x,y)\,\mathrm{d}S(y) + \int_\Omega G(x,y)f(y)\,\mathrm{d}y. \label{eq:green-sol}
\end{align}
\end{theorem}

\begin{remark}
Green 函数 $G(x,y)$ $(n=2,3)$ 的物理意义如下：在物体内部 $x$ 处放置一个单位点热源，与外界接触的表面保持零度，那么物体的稳定温度场就是 Green 函数；或者让某导体的表面接地，在其内部 $x$ 处放置一个单位正电荷，那么在导体内部所产生的电势分布也是 Green 函数。
\end{remark}

\begin{remark}
引入 Green 函数 $G(x,y)$ 的重要意义在于把求解具有任意非齐次项与任意边值的定解问题  归结为求解一个特定的边值问题 \eqref{eq:phi}。对于一些特殊区域，这样的特定的边值问题可以得到解的具体表达式。在一般情形，虽然不可能给出 Green 函数的表达式，但 Green 函数只依赖于区域，而与边值和非齐次项无关，这无论对理论研究还是对求解问题都带来很大的方便。
\end{remark}

\begin{theorem}[Green 函数的对称性]
对所有 $x,y\in\Omega$，$x\neq y$，有
\begin{align*}
G(y,x) = G(x,y).
\end{align*}
\end{theorem}

\begin{proof}
固定 $x,y\in\Omega$，$x\neq y$。取充分小 $\varepsilon>0$，使得 $B(x,\varepsilon)\cup B(y,\varepsilon)\subset\Omega$ 和 $B(x,\varepsilon)\cap B(y,\varepsilon)=\varnothing$。令
\begin{align*}
v(z) = G(x,z), \qquad w(z) = G(y,z),
\end{align*}
则
\begin{align*}
\Delta v(z) = 0,\quad (z\neq x), \qquad \Delta w(z) = 0,\quad (z\neq y).
\end{align*}
在区域 $\Omega\setminus [B(x,\varepsilon)\cup B(y,\varepsilon)]$ 上对 $v(z)$ 和 $w(z)$ 应用引理 \eqref{eq:green-second} 的公式得
\begin{align}
\int_{\partial B(x,\varepsilon)} \left( \frac{\partial v}{\partial n}w - \frac{\partial w}{\partial n}v \right) \mathrm{d}S(z) = \int_{\partial B(y,\varepsilon)} \left( \frac{\partial w}{\partial n}v - \frac{\partial v}{\partial n}w \right) \mathrm{d}S(z), \label{eq:sym-identity}
\end{align}
其中 $n$ 表示 $B(x,\varepsilon)\cup B(y,\varepsilon)$ 上单位内法向量。

我们首先考虑上式 \eqref{eq:sym-identity} 的左端项。由于 $\phi^x(z)$ 在 $\Omega$ 是有界的调和函数（见本章第三节）且 $w$ 在 $x$ 附近光滑，故
\begin{align*}
\left| \int_{\partial B(x,\varepsilon)} \frac{\partial w}{\partial n}v\,\mathrm{d}S(z) \right| \leqslant C\varepsilon^{n-1}\sup_{\partial B(x,\varepsilon)} |v|,
\end{align*}
其中 $C$ 是不依赖于 $\varepsilon$ 的常数。因此
\begin{align*}
\lim_{\varepsilon\to 0} \int_{\partial B(x,\varepsilon)} \frac{\partial v}{\partial n}v\,\mathrm{d}S(z) = 0.
\end{align*}
另外 $v(z)=\Gamma(z-x)-\phi^x(z)$，从而
\begin{align*}
\lim_{\varepsilon\to 0} \int_{\partial B(x,\varepsilon)} \frac{\partial v(z)}{\partial n}w(z)\,\mathrm{d}S(z)
&= \lim_{\varepsilon\to 0} \int_{\partial B(x,\varepsilon)} \frac{\partial}{\partial n}\Gamma(z-x)w(z)\,\mathrm{d}S(z) \\
&= \lim_{\varepsilon\to 0} \int_{\partial B(x,\varepsilon)} w(z)\,\mathrm{d}S(z) \\
&= w(x).
\end{align*}
因此，当 $\varepsilon\to 0$ 时，式 \eqref{eq:sym-identity} 的左端趋于 $w(x)$，而其右端趋于 $v(y)$。于是
\begin{align*}
w(x) = v(y),
\end{align*}
从而
\begin{align*}
G(y,x) = G(x,y).
\end{align*}
至此我们完成定理的证明。
\end{proof}

\subsubsection{半空间上的 Green 函数}

在以下的篇幅中，我们将构造出半空间 $\mathbb{R}_+^n$ 上的 Green 函数，并推导出半空间 $\mathbb{R}_+^n$ 上 Dirichlet 问题  的解的 Poisson 公式。

对于 $x=(x_1,\cdots,x_n)\in\mathbb{R}_+^n$，记它关于边界 $\partial\mathbb{R}_+^n=\mathbb{R}^{n-1}$ 的反射点为 $\tilde{x}=(x_1,\cdots,x_{n-1},-x_n)$。取
\begin{align*}
\phi^x(y) = \Gamma(y-\tilde{x}) \quad (x,y\in\mathbb{R}_+^n).
\end{align*}
注意到，当 $y\in\partial\mathbb{R}_+^n$ 时，$\phi^x(y)=\Gamma(y-x)$，于是 $\phi^x(y)$ 满足初值问题
\begin{align*}
\begin{cases}
-\Delta\phi^x(y) = 0, & y\in\mathbb{R}_+^n, \\
\phi^x = \Gamma(y-x), & y\in\partial\mathbb{R}_+^n.
\end{cases}
\end{align*}
因此 $\mathbb{R}_+^n$ 上的 Green 函数为
\begin{align*}
G(x,y) = \Gamma(y-x) - \phi^x(y) = \Gamma(y-x) - \Gamma(y-\tilde{x}).
\end{align*}

我们分别考虑下列两种情形：

(1) 当 $n\geqslant 3$ 时，由于
\begin{align*}
\frac{\partial}{\partial y_n}G(x,y) &= \frac{\partial}{\partial y_n}\Gamma(y-x) - \frac{\partial}{\partial y_n}\Gamma(y-\tilde{x}) \\
&= -\frac{1}{n\alpha(n)} \left( \frac{y_n-x_n}{|y-x|^n} - \frac{y_n+x_n}{|y-\tilde{x}|^n} \right),
\end{align*}
且边界 $\partial\mathbb{R}_+^n$ 的单位外法向量 $n=(0,\cdots,0,-1)$。因此，当 $x\in\mathbb{R}_+^n$，$y\in\partial\mathbb{R}_+^n$ 时，有
\begin{align}
\frac{\partial}{\partial n}G(x,y) = -\frac{\partial}{\partial y_n}G(x,y) = -\frac{2x_n}{n\alpha(n)|y-x|^n}, \label{eq:normal-deriv-half}
\end{align}
其中
\begin{align*}
|y-x|^2 = (y_1-x_1)^2 + \cdots + (y_{n-1}-x_{n-1})^2 + x_n^2.
\end{align*}

(2) 当 $n=2$ 时，由于
\begin{align*}
\frac{\partial}{\partial y_2}G(x,y) = \frac{\partial}{\partial y_2}\Gamma(y-x) - \frac{\partial}{\partial y_2}\Gamma(y-\tilde{x})
= -\frac{1}{2\pi}\left( \frac{y_2-x_2}{|y-x|^2} - \frac{y_2+x_2}{|y-\tilde{x}|^2} \right),
\end{align*}
且 $\alpha(2)=\pi$，因此，当 $x\in\mathbb{R}_+^2$，$y\in\partial\mathbb{R}_+^2$ 时，有
\begin{align*}
\frac{\partial}{\partial n}G(x,y) = -\frac{\partial}{\partial y_2}G(x,y) = -\frac{2x_2}{2\alpha(2)|y-x|^2},
\end{align*}
其中
\begin{align*}
|y-x|^2 = (y_1-x_1)^2 + x_2^2.
\end{align*}
综上所述，当 $n\geqslant 2$ 时，公式 \eqref{eq:normal-deriv-half} 成立。

假设 $u\in C^2(\mathbb{R}_+^n)\cap C(\overline{\mathbb{R}}_+^n)$ $(n\geqslant 2)$ 是方程
\begin{align}
\begin{cases}
-\Delta u = 0, & x\in\mathbb{R}_+^n, \\
u = g, & x\in\partial\mathbb{R}_+^n = \mathbb{R}^{n-1}
\end{cases} \label{eq:dirichlet-half}
\end{align}
的有界解。由定理 \eqref{eq:green-sol} 的解的表达式我们期望
\begin{align}
u(x) = \frac{2x_n}{n\alpha(n)}\int_{\mathbb{R}^{n-1}} \frac{g(y)}{|y-x|^n}\,\mathrm{d}y \label{eq:poisson-half}
\end{align}
是方程 \eqref{eq:dirichlet-half} 的解。通常称函数
\begin{align*}
K(x,y) = \frac{2x_n}{n\alpha(n)|y-x|^n} \quad (x\in\mathbb{R}_+^n,\; y\in\partial\mathbb{R}_+^n=\mathbb{R}^{n-1})
\end{align*}
为 $\mathbb{R}_+^n$ 的 Poisson 核，并称公式 \eqref{eq:poisson-half} 为 Poisson 公式。特别地，当 $n=2$ 时，记 $u(x)=u(x_1,x_2)$，则 Poisson 公式 \eqref{eq:poisson-half} 为
\begin{align*}
u(x_1,x_2) = \frac{x_2}{\pi}\int_{-\infty}^{+\infty} \frac{g(y)}{(y-x_1)^2+x_2^2}\,\mathrm{d}y.
\end{align*}

但是，上述 Poisson 公式 \eqref{eq:poisson-half} 只给出边值问题 \eqref{eq:dirichlet-half} 的形式解，因为此时区域 $\mathbb{R}_+^n$ 是无界的，公式 \eqref{eq:poisson-half} 的推导并不严格。为此我们需要验证公式 \eqref{eq:poisson-half} 的确给出边值问题 \eqref{eq:dirichlet-half} 的解。实际上，我们利用 Poisson 公式 \eqref{eq:poisson-half} 构造出在上半空间 $\mathbb{R}_+^n$ 的边界 $\partial\mathbb{R}_+^n$ 上指定取值的上半空间 $\mathbb{R}_+^n$ 上的调和函数。注意 Poisson 公式 \eqref{eq:poisson-half} 中 $u(x)$ 在上半空间的边界 $\partial\mathbb{R}_+^n$ 上并没有定义，因而在 Dirichlet 问题 \eqref{eq:dirichlet-half} 中，$u(x)$ 在边界 $\partial\mathbb{R}_+^n$ 上取边值 $g(x)$ 是在取极限的意义下成立的。

\begin{theorem}
假设 $g$ 是 $\mathbb{R}^{n-1}$ $(n\geqslant 2)$ 上有界连续函数。对于任意 $x\in\mathbb{R}_+^n$，$u(x)$ 由 Poisson 公式 \eqref{eq:poisson-half} 定义，即
\begin{align*}
u(x) = \int_{\mathbb{R}^{n-1}} K(x,y)g(y)\,\mathrm{d}y,
\end{align*}
则
\begin{enumerate}[(1)]
\item $u(x)$ 是 $\mathbb{R}_+^n$ 上无穷次可微的有界函数；
\item $\Delta u(x)=0,\ x\in\mathbb{R}_+^n$；
\item 对于任意 $x_0\in\partial\mathbb{R}_+^n$，当 $x\in\mathbb{R}_+^n$ 且 $x\to x_0$ 时，$u(x)\to g(x_0)$。
\end{enumerate}
\end{theorem}

\begin{proof}
(1) 对任意 $x\in\mathbb{R}_+^n$，我们可以证明
\begin{align}
\int_{\mathbb{R}^{n-1}} K(x,y)\,\mathrm{d}y = 1. \label{eq:poisson-half-normal}
\end{align}
由定理的假设，存在 $M>0$，使得 $|g|\leqslant M$。因此公式 \eqref{eq:poisson-half} 定义的 $u(x)$ 是有界的。当 $x\neq y$ 时，$K(x,y)$ 是光滑的。我们易证 $u\in C^\infty(\mathbb{R}_+^n)$。

(2) 当 $x\neq y$ 时，$\Delta_x G(x,y)=0$，于是 $\Delta_x \frac{\partial}{\partial y_n}G(x,y)=0$。所以，当 $x\in\mathbb{R}_+^n$，$y\in\partial\mathbb{R}_+^n$ 时，有 $\Delta_x K(x,y)=0$，且
\begin{align*}
\Delta u(x) = \int_{\mathbb{R}^{n-1}} \Delta_x K(x,y)g(y)\,\mathrm{d}y = 0.
\end{align*}

(3) 固定 $x_0\in\partial\mathbb{R}_+^n=\mathbb{R}^{n-1}$，$\varepsilon>0$。取 $\delta_1>0$ 足够小，使得当 $y\in\mathbb{R}^{n-1}$ 且 $|y-x_0|\leqslant \delta_1$ 时，有
\begin{align*}
|g(y)-g(x_0)| \leqslant \frac{\varepsilon}{2}.
\end{align*}
于是当 $x\in\mathbb{R}_+^n$ 且 $|x-x_0|\leqslant \frac{\delta_1}{2}$ 时，利用等式 \eqref{eq:poisson-half-normal} 我们计算得
\begin{align*}
|u(x)-g(x_0)| &= \left| \int_{\mathbb{R}^{n-1}} K(x,y)[g(y)-g(x_0)]\,\mathrm{d}y \right| \\
&\leqslant \int_{\mathbb{R}^{n-1}\cap B(x_0,\delta_1)} K(x,y)|g(y)-g(x_0)|\,\mathrm{d}y \\
&\quad + \int_{\mathbb{R}^{n-1}\setminus B(x_0,\delta_1)} K(x,y)|g(y)-g(x_0)|\,\mathrm{d}y \\
&\leqslant \frac{\varepsilon}{2} + 2M\int_{\mathbb{R}^{n-1}\setminus B(x_0,\delta_1)} K(x,y)\,\mathrm{d}y,
\end{align*}
这里 $B(x_0,\delta_1)$ 是 $\mathbb{R}^{n-1}$ 上以 $x_0$ 为心，$\delta_1$ 为半径的 $n-1$ 维球。

如果 $|x-x_0|\leqslant \frac{\delta_1}{2}$ 和 $|y-x_0|\geqslant \delta_1$，我们有
\begin{align*}
|x-x_0| \leqslant \frac{1}{2}|y-x_0|,
\end{align*}
从而
\begin{align*}
|y-x_0| \leqslant |y-x| + |x-x_0| \leqslant |y-x| + \frac{1}{2}|y-x_0|,
\end{align*}
因此
\begin{align*}
|y-x| \geqslant \frac{1}{2}|y-x_0|.
\end{align*}
当 $|x-x_0|\leqslant \frac{\delta_1}{2}$ 时，我们得到
\begin{align*}
\int_{\mathbb{R}^{n-1}\setminus B(x_0,\delta_1)} K(x,y)\,\mathrm{d}y
\leqslant x_n \frac{2^{n+1}}{n\alpha(n)} \int_{\mathbb{R}^{n-1}\setminus B(x_0,\delta_1)} |y-x_0|^{-n}\,\mathrm{d}y 
\leqslant C(n,\delta_1)x_n,
\end{align*}
这里 $C(n,\delta_1)$ 是仅依赖于 $n,\delta_1$ 的常数。取 $\delta_2>0$ 足够小，使当 $0<x_n\leqslant \delta_2$ 时，有
\begin{align*}
2M\int_{\mathbb{R}^{n-1}\setminus B(x_0,\delta_1)} K(x,y)\,\mathrm{d}y \leqslant \frac{\varepsilon}{2}.
\end{align*}
因此，当 $|x-x_0|\leqslant \delta = \min\{\frac{\delta_1}{2},\delta_2\}$ 时，有
\begin{align*}
|u(x)-g(x_0)| \leqslant \varepsilon.
\end{align*}
至此我们完成定理的证明。
\end{proof}

\subsubsection{球上的 Green 函数}

在以下的篇幅中，我们将构造出球 $B(0,R)$ 上的 Green 函数，并推导出球 $B(0,R)$ 上 Dirichlet 问题  的解的 Poisson 公式。

我们先构造出单位球 $B(0,1)$ 上的 Green 函数，然后通过伸缩变换获得一般球 $B(0,R)$ 上的 Green 函数。对于 $x\in\mathbb{R}^n\setminus\{0\}$，记它关于球面 $\partial B(0,1)$ 的对偶点为 $x^*=\frac{x}{|x|^2}$。显然 $|x^*||x|=1$。

我们将利用单位球的性质来构造其上的 Green 函数。首先我们需要求解如下问题：
\begin{align*}
\begin{cases}
-\Delta\phi^x(y) = 0, & y\in B(0,1), \\
\phi^x = \Gamma(y-x), & y\in\partial B(0,1).
\end{cases}
\end{align*}
我们分别考虑下列两种情形：

(1) 假设 $n\geqslant 3$。当 $y\neq x^*$ 时，$\Delta_y\Gamma(y-x^*)=0$，于是
\begin{align*}
\Delta_y\Gamma[|x|(y-x^*)] = 0.
\end{align*}
由对偶点 $x^*$ 与点 $x$ 的关系，当 $y\in\partial B(0,1)$ 时，$|x||y-x^*| = |y-x|$，于是当 $y\in\partial B(0,1)$ 时，$\Gamma[|x|(y-x^*)] = \Gamma(y-x)$。定义 $\phi^x(y)=\Gamma(|x|(y-x^*))$，则
\begin{align*}
\begin{cases}
-\Delta\phi^x(y) = 0, & y\in B(0,1), \\
\phi^x(y) = \Gamma(y-x), & y\in\partial B(0,1).
\end{cases}
\end{align*}
因此 $B(0,1)$ 上的 Green 函数为
\begin{align*}
G(x,y) = \Gamma(y-x) - \phi^x(y) = \Gamma(y-x) - \Gamma[|x|(y-x^*)].
\end{align*}
所以，当 $y\in\partial B(0,1)$ 时，有
\begin{align*}
\frac{\partial}{\partial y_i}G(x,y) &= \frac{\partial}{\partial y_i}\Gamma(y-x) - \frac{\partial}{\partial y_i}\Gamma[|x|(y-x^*)] \\
&= \frac{1}{n\alpha(n)}\left[ \frac{x_i-y_i}{|y-x|^n} + \frac{|x|^2(y_i-x_i^*)}{|x|^n|y-x^*|^n} \right] \\
&= \frac{1}{n\alpha(n)}\left( \frac{x_i-y_i}{|y-x|^n} + \frac{y_i|x|^2 - x_i}{|y-x|^n} \right) \\
&= -\frac{y_i(1-|x|^2)}{n\alpha(n)|y-x|^n},\quad i=1,2,\cdots,n.
\end{align*}
由于单位球 $B(0,1)$ 上的单位外法向量 $n=(y_1,y_2,\cdots,y_n)$ 且
\begin{align*}
\frac{\partial}{\partial n}G(x,y) = DG(x,y)\cdot n = \sum_{i=1}^n y_i\frac{\partial}{\partial y_i}G(x,y),
\end{align*}
故我们得到
\begin{align}
\frac{\partial}{\partial n}G(x,y) = -\frac{1-|x|^2}{n\alpha(n)|x-y|^n}. \label{eq:normal-deriv-ball}
\end{align}

(2) 假设 $n=2$。由基本解的表达式我们知道
\begin{align*}
G(x,y) = \Gamma(y-x) - \Gamma[|x|(y-x^*)]
\end{align*}
仍然是 $B(0,1)$ 上的 Green 函数。当 $y\in\partial B(0,1)$ 时，$|x||y-x^*|=|y-x|$。因此，当 $y\in\partial B(0,1)$ 时，有
\begin{align*}
\frac{\partial}{\partial y_i}G(x,y) &= \frac{\partial}{\partial y_i}\Gamma(y-x) - \frac{\partial}{\partial y_i}\Gamma[|x|(y-x^*)] \\
&= \frac{x_i-y_i}{2\pi|y-x|^2} + \frac{|x|^2(y_i-x_i^*)}{2\pi|x|^2|y-x^*|^2} \\
&= \frac{x_i-y_i}{2\pi|y-x|^2} + \frac{y_i|x|^2 - x_i}{2\pi|y-x|^2} \\
&= -\frac{y_i(1-|x|^2)}{2\alpha(2)|y-x|^2},\quad i=1,2,
\end{align*}
从而
\begin{align*}
\frac{\partial}{\partial n}G(x,y) = \sum_{i=1}^2 y_i\frac{\partial}{\partial y_i}G(x,y) = -\frac{1-|x|^2}{2\alpha(2)|x-y|^2}.
\end{align*}
综上所述，当 $n\geqslant 2$ 时，公式 \eqref{eq:normal-deriv-ball} 成立。

假设 $u\in C^2(B(0,1))\cap C(\overline{B}(0,1))$ 是
\begin{align}
\begin{cases}
-\Delta u = 0, & x\in B(0,1), \\
u = g, & x\in\partial B(0,1)
\end{cases} \label{eq:dirichlet-unit}
\end{align}
的解，则由定理 \eqref{eq:green-sol} 的表达式我们得到
\begin{align*}
u(x) = -\int_{\partial B(0,1)} g(y)\frac{\partial}{\partial n}G(x,y)\,\mathrm{d}S(y).
\end{align*}
于是由公式 \eqref{eq:normal-deriv-ball} 得到
\begin{align*}
u(x) = \frac{1-|x|^2}{n\alpha(n)}\int_{\partial B(0,1)} \frac{g(y)}{|x-y|^n}\,\mathrm{d}S(y).
\end{align*}

假设 $u\in C^2(B(0,R))\cap C(\overline{B}(0,R))$ 是
\begin{align}
\begin{cases}
-\Delta u = 0, & x\in B(0,R), \\
u = g, & x\in\partial B(0,R)
\end{cases} \label{eq:dirichlet-ball}
\end{align}
的解，则 $\tilde{u}(x)=u(Rx)$ 满足 Dirichlet 问题 \eqref{eq:dirichlet-unit}，其中用 $\tilde{g}(x)=g(Rx)$ 代替 $g$。我们作变量替换得到
\begin{align}
u(x) = \frac{R^2-|x|^2}{n\alpha(n)R} \int_{\partial B(0,R)} \frac{g(y)}{|x-y|^n}\,\mathrm{d}S(y). \label{eq:poisson-ball}
\end{align}
通常称函数
\begin{align*}
K(x,y) = \frac{R^2-|x|^2}{n\alpha(n)R|x-y|^n} \quad (x\in B(0,R),\; y\in\partial B(0,R))
\end{align*}
为 Poisson 核，而称公式 \eqref{eq:poisson-ball} 为 Poisson 公式。实际上，我们利用 Poisson 公式 \eqref{eq:poisson-ball} 构造出在球面 $\partial B(0,R)$ 上指定取值的球 $B(0,R)$ 上的调和函数。

以下我们验证公式 \eqref{eq:poisson-ball} 的确给出 Dirichlet 问题 \eqref{eq:dirichlet-ball} 的解。注意公式 \eqref{eq:poisson-ball} 中 $u(x)$ 在球面 $\partial B(0,R)$ 上没有定义，因而在 Dirichlet 问题 \eqref{eq:dirichlet-ball} 中，$u(x)$ 在球面上取边值 $g(x)$ 是在取极限的意义下成立的。

\begin{theorem}
假设 $B(0,R)\subset\mathbb{R}^n$，$g:\partial B(0,R)\to\mathbb{R}$ 连续，对于任意 $x\in B(0,R)$，函数 $u(x)$ 由 Poisson 公式 \eqref{eq:poisson-ball} 定义，即
\begin{align*}
u(x) = \int_{\partial B(0,R)} K(x,y)g(y)\,\mathrm{d}S(y),
\end{align*}
则
\begin{enumerate}[(1)]
\item $u(x)$ 是 $B(0,R)$ 上无穷次可微的有界函数；
\item $\Delta u(x)=0,\ x\in B(0,R)$；
\item 任给 $x_0\in\partial B(0,R)$，当 $x\in B(0,R)$ 且 $x\to x_0$ 时，$u(x)\to g(x_0)$。
\end{enumerate}
\end{theorem}
\begin{remark}
当 $n=2$ 时，在公式 \eqref{eq:poisson-ball} 中用极坐标表示 $x=(\rho\cos\theta,\rho\sin\theta)$，$y=(R\cos\varphi,R\sin\varphi)$，记 $u(x)=u(\rho,\theta)$，则公式 \eqref{eq:poisson-ball} 可以表示成
\begin{align*}
u(\rho,\theta) = \frac{1}{2\pi}\int_0^{2\pi} \frac{R^2-\rho^2}{R^2+\rho^2-2R\rho\cos(\varphi-\theta)}g(\varphi)\,\mathrm{d}\varphi,
\end{align*}
其中 $g(\varphi)=g(R\cos\varphi,R\sin\varphi)$。实际上，我们在“复变函数”这门课程中学过以上公式。
\end{remark}
\begin{proof}
(1) 对任意 $x\in B(0,R)$，我们对特殊情形 $u=1$ 应用公式 \eqref{eq:green-sol} 得到
\begin{align*}
\int_{\partial B(0,R)} K(x,y)\,\mathrm{d}S(y) = 1.
\end{align*}
由定理的假设，存在 $M>0$，使得 $|g|\leqslant M$。因此公式 \eqref{eq:poisson-ball} 定义的 $u(x)$ 是有界的。当 $x\neq y$ 时，$K(x,y)$ 是光滑的，我们易证 $u\in C^\infty(B(0,R))$。

(2) 当 $x\neq y$ 时，$\Delta_x G(x,y)=0$，于是 $\Delta_x \frac{\partial}{\partial n}G(x,y)=0$。从而，当 $x\in B(0,R)$，$y\in\partial B(0,R)$ 时，$\Delta_x K(x,y)=0$ 且
\begin{align*}
\Delta u(x) = \int_{\partial B(0,R)} \Delta_x K(x,y)g(y)\,\mathrm{d}S(y) = 0.
\end{align*}

(3) 固定 $x_0\in\partial B(0,R)$，$\varepsilon>0$。取 $\delta_1>0$ 足够小，使得当 $y\in\partial B(0,R)$ 且 $|y-x_0|\leqslant \delta_1$ 时，
\begin{align*}
|g(y)-g(x_0)| \leqslant \frac{\varepsilon}{2},
\end{align*}
于是当 $x\in B(0,R)$ 且 $|x-x_0|\leqslant \frac{\delta_1}{2}$ 时，利用归一化等式我们计算得
\begin{align*}
|u(x)-g(x_0)| &= \left| \int_{\partial B(0,R)} K(x,y)[g(y)-g(x_0)]\,\mathrm{d}S(y) \right| \\
&\leqslant \int_{\partial B(0,R)\cap\{|y-x_0|\leqslant \delta_1\}} K(x,y)|g(y)-g(x_0)|\,\mathrm{d}S(y) \\
&\quad + \int_{\partial B(0,R)\cap\{|y-x_0|>\delta_1\}} K(x,y)|g(y)-g(x_0)|\,\mathrm{d}S(y) \\
&\leqslant \frac{\varepsilon}{2} + \frac{2M(R^2-|x|^2)R^{n-2}}{(\delta_1/2)^n}.
\end{align*}
取 $\delta>0$ 足够小，当 $|x-x_0|<\delta$ 时，$R^2-|x|^2$ 足够小，于是我们得到
\begin{align*}
|u(x)-g(x_0)| \leqslant \varepsilon.
\end{align*}
至此我们完成定理的证明。
\end{proof}







\end{document}